\documentclass[12 pt, twoside]{amsart}

\usepackage{../current_style} % Must be in a child of recipes

\begin{document}

\title{Cardamom Rusks}
\maketitle
\subtitle{30 minutes plus 30 more of chilling and almost 90 for baking}

\noindent \desc{These tea cookies, ``rusks'', originally for dipping in a milk-like soup work fantastically for dipping in tea. Originally taken from Nicole Accettola's ``Scandinavian from Scratch'', I've modified them to include more cardamom, a bit of nutmeg, and some time to chill in the fridge. I've also changed the baking process a decent bit, cooking at a little bit higher of a temperature, for a bit longer, and flipping to get a nicer and more uniform color on each side.}

\recipe{
\begin{ingredients}
\item 1 stick of butter
\item $\frac{1}{4}$ cup sugar
\item 2 tbsp whole milk
\item 2 eggs
\item 2 gups all purpose flour
\item $\frac{3}{4}$ tsp baking powder
\item $\frac{1}{4}$ tsp salt
\item 2 tsp cardamom
\item 1 tsp nutmeg
\end{ingredients}
}{
\begin{directions}
\item In a bowl (preferably for a stand mixer) cream the butter and sugar until the sugar has disappeared. Add milk and eggs and beat some more---they should be pretty well incorporated, but don't worry if they aren't.
\item In a separate bowl combine the flour, baking powder, salt, cardamom, and nutmeg. Add to the butter mixture and mix together for a few minutes, until there's a somewhat dry dough---you may need to knead by hand. Let rest in the fridge for a half hour.
\item Preheat the oven to 325. Divide dough into four pieces and then roll each into a log. I roll them so they have a bout an inch of diameter, but if you want a larger or smaller cookie you can roll them more or less.
\item Bake for about 15 minutes, until just starting to get some color. Lower the heat to 225 and cut the logs into roughly cm inch slices. 
\item Place face down on the baking sheet and bake for 30-40 minutes, or until the face-down side is golden. Then flip them over and bake for another 15-20 minutes.
\end{directions}
}

\end{document}
